% ============================================================
% f2014_hunan_science_q12_circle.tex
% 2014 湖南卷（理科）第 12 题：圆 O 的两条弦 AB、BC
% 来源 PDF：普通高考/2014/2014湖南理.pdf
% 原页：第 1 页（A3 横排）第 2 栏中部，印面 885，第 12 题题干下方
% 结构：O 为圆心，半径 R=3/2；A(-3/2,0) 在圆上，AO 为半径；
%       BC 为垂直于 AO 的弦，B、C 分别在 AO 上下，BM=MC=√2，
%       OM=1/2。比例 1.3 cm/单位。
% ============================================================
\begin{tikzpicture}[>=Stealth, line width=0.7pt, font=\small,
  x=1.3cm, y=1.3cm, line cap=round, line join=round]
  % 圆
  \draw (0,0) circle (1.5);
  % 半径 AO
  \draw (-1.5,0) -- (0,0);
  % 弦 BC（过 AO 上的点 M）
  \draw (-0.5,1.41421) -- (-0.5,-1.41421);
  % 弦 AB
  \draw (-1.5,0) -- (-0.5,1.41421);
  % 圆心与顶点
  \fill (0,0) circle (1.4pt);
  \node[right] at (0.06,0.02) {$O$};
  \node[left]  at (-1.56,0.10) {$A$};
  \node[above] at (-0.52,1.52) {$B$};
  \node[below] at (-0.52,-1.52) {$C$};
\end{tikzpicture}
